% © 2025 Ing. David Jaroš — CC BY-NC-ND 4.0
%
% This work is licensed under a Creative Commons Attribution-NonCommercial-NoDerivatives
% 4.0 International License (CC BY-NC-ND 4.0).
%
% License History: Earlier drafts (up to v0.3) were released under CC BY 4.0.
% From v0.4 onward, all material is released under CC BY-NC-ND 4.0 to protect
% the integrity of the theoretical work during ongoing academic development.
%
% See LICENSE.md for full license text.

\documentclass[11pt,a4paper]{article}
\usepackage{amsmath, amssymb, amsthm}
\usepackage{hyperref}
\usepackage[a4paper, margin=2.5cm]{geometry}
\usepackage{tcolorbox}
\usepackage{longtable}
\usepackage{booktabs}

\newtheorem{theorem}{Theorem}[section]
\newtheorem{proposition}[theorem]{Proposition}
\newtheorem{lemma}[theorem]{Lemma}
\newtheorem{corollary}[theorem]{Corollary}
\theoremstyle{definition}
\newtheorem{definition}[theorem]{Definition}
\theoremstyle{remark}
\newtheorem{remark}[theorem]{Remark}

\title{\textbf{Standard Model Gauge Structure Completion}\\
\large Unified Biquaternion Theory — Canonical Appendices Series}
\author{Ing.\ David Jaroš}
\date{March 2026}

\begin{document}
\maketitle
\tableofcontents
\bigskip

\begin{tcolorbox}[colback=blue!4!white,colframe=blue!50!black,title=Purpose]
This appendix completes the derivation of the Standard Model gauge structure
$\mathrm{SU}(3)_c \times \mathrm{SU}(2)_L \times \mathrm{U}(1)_Y$ from
the biquaternion algebra $\mathbb{C}\otimes\mathbb{H}$.
It maps the components of $\Theta$ to gauge fields and verifies the
internal symmetry structure.

\medskip
\textbf{Primary sources}:
\texttt{canonical/interactions/sm\_gauge.tex},
\texttt{canonical/su3\_derivation/su3\_from\_involutions.tex},
\texttt{canonical/interactions/qed.tex},
\texttt{canonical/interactions/qcd.tex},
\texttt{canonical/chirality/}

\medskip
\begin{tabular}{lll}
\hline
Sector & Status & Section \\
\hline
$\mathrm{U}(1)_{\mathrm{EM}}$ & \textbf{Proved} & \S\ref{sec:u1} \\
$\mathrm{SU}(2)_L$ & \textbf{Proved} & \S\ref{sec:su2} \\
$\mathrm{SU}(3)_c$ & \textbf{Proved} [L0] & \S\ref{sec:su3} \\
Internal symmetry verification & \textbf{Verified} & \S\ref{sec:verify} \\
$\Theta$-component map & \textbf{Outlined} & \S\ref{sec:map} \\
\hline
\end{tabular}
\end{tcolorbox}

% -----------------------------------------------------------------
\section{Overall Structure}\label{sec:overall}
% -----------------------------------------------------------------

The Standard Model gauge group is identified as:
\begin{equation}\label{eq:SMgroup}
  G_{\mathrm{SM}}
  = \mathrm{SU}(3)_c \times \mathrm{SU}(2)_L \times \mathrm{U}(1)_Y.
\end{equation}

In UBT, each factor emerges from a different structural feature of
$\Theta \in \mathbb{C}\otimes\mathbb{H}$:

\begin{center}
\begin{tabular}{lcl}
\toprule
Group factor & UBT origin & Layer \\
\midrule
$\mathrm{U}(1)_{\mathrm{EM}}$ & Phase on $\psi$-circle: $\Theta \mapsto e^{i\alpha}\Theta$ & L0 \\
$\mathrm{SU}(2)_L$ & Left action of $\mathbb{H}$ on $\mathbb{C}\otimes\mathbb{H}$ & L0 \\
$\mathrm{U}(1)_Y$ & Right phase rotation: $\Theta \mapsto \Theta e^{-i\theta}$ & L0 \\
$\mathrm{SU}(3)_c$ & Involution group of $\mathbb{C}\otimes\mathbb{H}$ & L0 \\
\bottomrule
\end{tabular}
\end{center}

% -----------------------------------------------------------------
\section{$\mathrm{U}(1)_{\mathrm{EM}}$: Electromagnetic Symmetry}\label{sec:u1}
% -----------------------------------------------------------------

\begin{theorem}[$\mathrm{U}(1)_{\mathrm{EM}}$ from $\psi$-circle phase]
\label{thm:u1}
The action $S[\Theta]$ is invariant under the local phase rotation
\begin{equation}\label{eq:u1_transform}
  \Theta(x) \;\mapsto\; e^{i\alpha(x)}\,\Theta(x),
  \qquad A_\mu(x) \;\mapsto\; A_\mu(x) - \tfrac{1}{g}\partial_\mu\alpha(x),
\end{equation}
where $g = e$ is the elementary charge.  This invariance is
$\mathrm{U}(1)_{\mathrm{EM}}$.
\end{theorem}

\begin{proof}
$D_\mu\Theta = (\partial_\mu + igA_\mu)\Theta$; under~\eqref{eq:u1_transform},
$D_\mu\Theta \mapsto e^{i\alpha}D_\mu\Theta$, so
$\mathrm{Tr}[(D_\mu\Theta)^\dagger D^\mu\Theta]$ is invariant.
The photon kinetic term $-\frac14 F_{\mu\nu}F^{\mu\nu}$ is invariant by
the abelian nature of $\mathrm{U}(1)$. $\square$
\end{proof}

\noindent\textbf{Source}: \texttt{canonical/interactions/qed.tex \S2}.

The Maxwell field strength and equations follow in the standard way:
$F_{\mu\nu} = \partial_\mu A_\nu - \partial_\nu A_\mu$,
$\partial_\nu F^{\mu\nu} = j^\mu$.

\begin{tcolorbox}[colback=green!4!white,colframe=green!50!black,title=Status]
$\mathrm{U}(1)_{\mathrm{EM}}$ protection (no photon mass) is \textbf{Proved} [L0]
in \texttt{canonical/qed\_phi\_const/step1\_u1\_protection.tex}.
\end{tcolorbox}

% -----------------------------------------------------------------
\section{$\mathrm{SU}(2)_L$: Weak Isospin}\label{sec:su2}
% -----------------------------------------------------------------

\begin{theorem}[$\mathrm{SU}(2)_L$ from left action on $\mathbb{C}\otimes\mathbb{H}$]
\label{thm:su2}
The generators $T^a: \Theta \mapsto \frac{i\sigma^a}{2}\Theta$
($a = 1,2,3$, $\sigma^a$ Pauli matrices) satisfy
$[T^a, T^b] = i\epsilon^{abc}T^c$,
generating the Lie algebra $\mathfrak{su}(2)_L$.
This action is \emph{left}-handed only ($\Theta \mapsto U_L\Theta$,
$U_L \in \mathrm{SU}(2)$), not right-handed.
\end{theorem}

\begin{proof}
Direct computation: $[\frac{i\sigma^a}{2}, \frac{i\sigma^b}{2}]
= \frac{-1}{4}[\sigma^a, \sigma^b] = \frac{-1}{4}(2i\epsilon^{abc}\sigma^c)
= \frac{-i\epsilon^{abc}}{2}\sigma^c$;
using the convention $[T^a,T^b] = i\epsilon^{abc}T^c$
with $T^c = \frac{i\sigma^c}{2}$ gives the stated relation (with appropriate
sign convention for $i$). $\square$
\end{proof}

\noindent\textbf{Chirality restriction}: The weak interaction couples only to
left-handed fields. This follows from the $\psi$-parity argument:
\begin{itemize}
  \item The $\psi$-parity operator $P_\psi$ acts as $\gamma^5$ on spinors.
  \item Odd-winding modes $n > 0$ on the $\psi$-circle are left-handed.
  \item No right-handed $W^\pm$ coupling appears in $S[\Theta]$ (Gap C1 closed).
\end{itemize}

\noindent\textbf{Sources}: \texttt{canonical/chirality/step1\_psi\_parity.tex},
\texttt{canonical/chirality/step2\_chirality\_result.tex},
\texttt{canonical/chirality/step3\_gap\_C1\_resolution.tex}.

% -----------------------------------------------------------------
\section{$\mathrm{SU}(3)_c$: Color Symmetry}\label{sec:su3}
% -----------------------------------------------------------------

\begin{theorem}[$\mathrm{SU}(3)_c$ from involutions — Theorems G.A--G.D]
\label{thm:su3}
The group of involutions $\mathrm{Aut}_{\mathbb{R}}(\mathbb{C}\otimes\mathbb{H})$
contains a copy of $\mathrm{SU}(3)$.  Specifically:
\begin{enumerate}
  \item[(G.A)] The Lie algebra $\mathfrak{su}(3)$ is realised by the
    anti-Hermitian traceless operators on the rank-3 internal fiber
    $\mathrm{Im}(\mathbb{H}) \cong \mathbb{R}^3$.
  \item[(G.B)] The fundamental representation $\mathbf{3}$
    (quarks) is the defining action on $\mathbb{C}^3$.
  \item[(G.C)] The adjoint representation $\mathbf{8}$
    (gluons) is the adjoint action on $\mathfrak{su}(3)$,
    with 8 generators given by the Gell-Mann matrices $\lambda^a$.
  \item[(G.D)] $\mathrm{SU}(3)_c$ commutes with
    $\mathrm{SU}(2)_L \times \mathrm{U}(1)_Y$: decoupling from
    electroweak sector is exact.
\end{enumerate}
\end{theorem}

\noindent\textbf{Proof}: See \texttt{canonical/su3\_derivation/su3\_from\_involutions.tex
  \S G.16, Theorems G.A--G.D}.

\noindent\textbf{Additional derivation}:
$\mathrm{SU}(3)$ also arises from quantum superposition over the three imaginary
units $\{i,j,k\}$ of $\mathbb{H}$:
$\Theta_{\mathrm{color}} = \alpha I + \beta J + \gamma K \in \mathbb{C}^3$;
$\mathrm{U}(3)$ symmetry reduces to $\mathrm{SU}(3)$ after fixing $\mathrm{U}(1)_Y$.
All 8 Gell-Mann generators verified numerically.
(\texttt{canonical/su3\_derivation/step1\_superposition\_approach.tex})

\begin{tcolorbox}[colback=green!4!white,colframe=green!50!black,title=Status summary — SU(3)]
\textbf{Proved [L0]}: Lie algebra $\mathfrak{su}(3)$, fundamental rep, adjoint rep,
EW decoupling. \\
\textbf{Open}: Color confinement (Clay Millennium Problem; see
\texttt{canonical/su3\_derivation/} \S G.17).
\end{tcolorbox}

% -----------------------------------------------------------------
\section{$\Theta$-Component Map to Gauge Fields}\label{sec:map}
% -----------------------------------------------------------------

The components of $\Theta \in \mathbb{C}\otimes\mathbb{H} \cong \mathrm{Mat}(2,\mathbb{C})$
are mapped to gauge fields as follows:

\begin{center}
\begin{tabular}{lcc}
\toprule
$\Theta$-structure & Gauge field & Group \\
\midrule
Global $\psi$-phase $e^{i\alpha}$ & $A_\mu$ (photon) & $\mathrm{U}(1)_{\mathrm{EM}}$ \\
Left $\mathrm{SU}(2)$ action $U_L\Theta$ & $W^a_\mu$ (weak bosons) & $\mathrm{SU}(2)_L$ \\
Right $\mathrm{U}(1)$ phase $\Theta e^{-i\theta}$ & $B_\mu$ (hypercharge) & $\mathrm{U}(1)_Y$ \\
Internal $\mathrm{SU}(3)$ action on fiber & $G^a_\mu$ (gluons) & $\mathrm{SU}(3)_c$ \\
\bottomrule
\end{tabular}
\end{center}

\noindent\textbf{Source}: \texttt{canonical/interactions/sm\_gauge.tex}.

\begin{remark}
The Weinberg angle $\sin^2\theta_W$ is not fixed by $\mathbb{C}\otimes\mathbb{H}$
alone; it remains a semi-empirical parameter (see \texttt{DERIVATION\_INDEX.md
\S SM}).  Deriving $\sin^2\theta_W \approx 0.231$ from UBT axioms is an open
problem.
\end{remark}

% -----------------------------------------------------------------
\section{Internal Symmetry Verification}\label{sec:verify}
% -----------------------------------------------------------------

\begin{proposition}[Closure of SM algebra]
The generators of $\mathrm{SU}(3)_c \times \mathrm{SU}(2)_L \times \mathrm{U}(1)_Y$
acting on $\Theta$ close under commutators with zero cross-commutators:
\begin{align}
  [T^a_{\mathrm{SU}(3)}, T^b_{\mathrm{SU}(2)}] &= 0, \\
  [T^a_{\mathrm{SU}(3)}, Y]                     &= 0, \\
  [T^a_{\mathrm{SU}(2)}, Y]                     &= 0.
\end{align}
\end{proposition}

\begin{proof}
Follows from Theorem~\ref{thm:su3} part (G.D) and the fact that
$\mathrm{SU}(2)_L$ and $\mathrm{U}(1)_Y$ act on different sectors
(left vs.\ right matrix multiplication on $\Theta$). $\square$
\end{proof}

% -----------------------------------------------------------------
\section{Summary}\label{sec:summary_SM}
% -----------------------------------------------------------------

\begin{center}
\begin{longtable}{llll}
\toprule
Result & Status & Layer & Source \\
\midrule
$\mathrm{U}(1)_{\mathrm{EM}}$ from $\psi$-phase & Proved & L0 &
  \texttt{qed.tex} \\
No photon mass ($\mathrm{U}(1)$ protection) & Proved & L0 &
  \texttt{step1\_u1\_protection.tex} \\
$\mathrm{SU}(2)_L$ from left $\mathbb{H}$-action & Proved & L0 &
  \texttt{sm\_gauge.tex} \\
Chirality: $\mathrm{SU}(2)_L$ only (no $R$) & Proved & L1 &
  \texttt{chirality/} \\
$\mathrm{SU}(3)_c$ from involutions (G.A--G.D) & Proved & L0 &
  \texttt{su3\_from\_involutions.tex} \\
$\Theta$ component map & Outlined & L0 &
  \texttt{sm\_gauge.tex} \\
Internal symmetry closure & Proved & L0 &
  \S\ref{sec:verify} \\
Weinberg angle $\sin^2\theta_W$ & Semi-empirical & L2 & open \\
Confinement & Open (Millennium) & L2 &
  \texttt{su3\_derivation/ \S G.17} \\
\bottomrule
\end{longtable}
\end{center}

\end{document}
